\clearpage
\fancyhead[L]{\small CONNAISSANCE DES TEMS.}
\SourceLeaf{614}
\begin{center}\small Göttingische gelehrte Anzeigen. Stück 40 und 41. Seiten 402..408. 11. März 1819.\\[3pt]\rule{.72\linewidth}{.3pt}\end{center}

\noindent\emph{Connaissance des tems ou des mouvemens célestes à l'usage des astronomes et des navigateurs pour l'an {\upshape 1820}, publiée par le bureau des longitudes.} Paris 1818. Bei der Witwe Courcier. 446 Seiten. Großoktav.

\begingroup\brokenpenalty=10000\relax
In der Einrichtung des Kalenders finden wir in diesem Jahrgang nur eine Änderung. Die Namen der Heiligen, die sonderbarerweise seit ein paar Jahren zwischen die astronomischen Rubriken gesetzt worden waren, sind wieder weggelassen worden. Den dadurch gewonnenen Platz hat man dazu verwendet, beim Abstand des Widderpunkts von der Sonne und bei der Deklination der Sonne eine Differenzspalte hinzuzufügen. Stattdessen hätte man wohl etwas Brauchbareres wählen können. Bei der Tafel zur geografischen Lage der wichtigsten Orte der Erde wird versichert, dass sie erneut durchgesehen worden sei und dass man diese neueste Ausgabe als die genaueste und zuverlässigste anzusehen habe. Diese Versicherung ist jedoch ein stehender Bestandteil, und man braucht in diesem Band ebenso wie in seinen Vorgängern nicht lange zu suchen, um sich davon zu überzeugen, dass dieses Verzeichnis sehr nachlässig redigiert wird. So finden wir z.\,B., die Polhöhe von Ofen um $32''$, die von Göttingen um $24' 26''$ zu groß, die von Königsberg um $38''$ zu klein und die von Mannheim um $5''$ zu groß. Beim ersten und dritten dieser Orte mögen andere Punkte als die gegenwärtigen Sternwarten zugrunde liegen. Das wenigstens hätte jedoch nicht der Fall sein sollen. Von den Beigaben gehen wir hier nur auf die Originalaufsätze ein und übergehen stillschweigend die auch in diesem Band sehr zahlreichen, weitschweifigen Auszüge aus anderen Büchern und Abhandlungen. Den Anfang machen die wertvollen Tafeln für Vesta von \textsc{Daussy}, die auf den im Jahrgang 1818 abgedruckten Elementen und Störungsformeln beruhen. Darauf folgt eine lehrreiche Abhandlung von \textsc{Laplace} über den Apparat zur Bestimmung der Länge des Sekundenpendels. Darin wird besonders der Einfluss des Umstands untersucht, dass die Schneide, auf der das Pendel schwingt, streng genommen keine mathematische Linie, sondern ein Stück einer kleinen Zylinderfläche ist. Die vortrefflichen Versuche von \textsc{Borda} zu diesem Gegenstand verdienten diese subtile Berücksichtigung. --- Über die Bestimmung der Senkrechten zum Erdmeridian und über verschiedene andere damit zusammenhängende Aufgaben, von \textsc{Puissant}. Unseres Erachtens ist die gesamte bisherige Behandlungsweise dieser Gegenstände noch nicht die richtige. Erst wenn man sie aus neuen Blickwinkeln betrachtet (was hier nicht geschehen ist), darf man hoffen, Vollständigkeit und Genauigkeit mit Einfachheit und Flexibilität bei den Rechenoperationen zu verbinden. Es scheint umso wünschenswerter, dass die Geometer ihre Aufmerksamkeit darauf richten, da gegenwärtig mehrere Regierungen groß angelegte Vermessungen veranlassen, die, mit den besten Hilfsmitteln ausgeführt, auf neue und wichtige Erkenntnisse über die Gestalt der Erde hoffen lassen. Nach einer hier von Herrn \textsc{Puissant} gegebenen Mitteilung hat auch die französische Regierung eine neue Triangulation beschlossen, die an die Meridianlinie von Dünkirchen und an die Senkrechte von Brest nach Straßburg angeschlossen und auf das ganze Königreich ausgedehnt werden soll. --- Über die elliptische Bahn des Kometen von 1783 und ihre Ähnlichkeit mit der Bahn des Kometen von 1793, von \textsc{Burckhardt}. Die geringe Übereinstimmung der Beobachtungen des Kometen von 1783 mit den von verschiedenen Astronomen erprobten parabolischen Elementen veranlasste Herrn \textsc{Burckhardt}, diese Bahn als elliptisch zu berechnen. Herr \textsc{Burckhardt} erhielt eine Ellipse mit einer Umlaufzeit von $5\tfrac{1}{2}$ Jahren, die die Beobachtungen bei weitem besser darstellt. Auch eine Ellipse mit einer Umlaufzeit von 10 Jahren ergab noch eine recht gute Übereinstimmung. Da diese Elemente einige Ähnlichkeit mit denen des zweiten Kometen von 1793 haben, schien es denkbar, dass beide Kometen identisch wären. Deshalb unterzog \SourceLeaf{615}Herr \textsc{Burckhardt} auch die Beobachtungen des Kometen von 1793, die bisher noch nicht gedruckt worden waren, einer neuen Berechnung. Leider sind diese Beobachtungen wenig genau. Es ist kaum möglich, anhand der hier mitgeteilten Ergebnisse ein bestimmtes Urteil zu fällen, zumal sich dabei wahrscheinlich Druckfehler eingeschlichen haben, deren Berichtigung sich nicht erraten lässt (die Länge des Knotens, die \textsc{Saron} für eine parabolische Bahn als $83^\circ\,55'$ ermittelt hatte, wird von Herrn \textsc{Burckhardt} für die eine elliptische Bahn mit $359'\,4'\,48''$ (sic) angegeben und bei einer zweiten elliptischen Bahn ganz ausgelassen). Gleichwohl verdient Herr \textsc{Burckhardt} Dank dafür, dass er die Aufmerksamkeit der Astronomen auf diesen Kometen gelenkt und die Beobachtungen von \textsc{Messier} zu diesem Kometen hier erstmals bekannt gemacht hat. --- Auszug aus einem Schreiben von Herrn \textsc{Schumacher} an Herrn \textsc{Delambre}. Die Unterschiede zwischen \textsc{Tycho}s und \textsc{Picard}s Azimutbestimmungen auf der Insel Hueen werden hier anhand einer Abhandlung \textsc{Augustin}s befriedigend aufgeklärt. Der Turm von Helsingborg, den \textsc{Picard} beobachtet hatte, war der falsche. Er war erst nach \textsc{Tycho}s Zeit erbaut worden. Außerdem hatte \textsc{Tycho} anfangs mit unvollkommenen Instrumenten seine Mittagslinien um einen Viertelgrad falsch gezogen. Diesen Fehler hatte er aber laut der Historia Coelestis schon einige Jahre später selbst bemerkt und berichtigt. --- Über einen Aufsatz von Herrn \textsc{Littrow} über die thermometrische Korrektur der Strahlenbrechung. Herr \textsc{Delambre} verteidigt hier sehr ausführlich seine eigene Tafel, der Herr \textsc{Littrow} Ungenauigkeit vorgeworfen hatte. Doch obwohl diese Rechtfertigung 10 Seiten einnimmt, wird nicht klar hervorgehoben, worin eigentlich der Fehler von Herrn \textsc{Littrow} bestand. Herr \textsc{Littrow} hatte voreilig angenommen, wenn die Dichte der Luft, die beim Gefrierpunkt als Einheit angesetzt wird, für eine Temperatur von $\alpha$ Grad darüber durch die Formel
\[
\frac{1}{1+0.0046875\,\alpha}
\]
 ausgedrückt wird, bleibe \emph{dieselbe} Formel auch dann gültig, wenn die Dichte bei $8^\circ$ als Einheit angesetzt wird und $\alpha$ dann den Thermometerstand über $8^\circ$ bezeichnet. Diese Annahme ist offenbar unrichtig. --- Über ein neues Mittel, die Dauer der Pendelschwingungen zu regulieren, von Herrn \textsc{Prony}. Die Pendelstange ist etwas nach oben verlängert. An ihr ist eine sehr dünne, kurze Querstange durch Reibung so befestigt, dass sie sich in einer zur Pendelstange senkrechten Ebene drehen lässt. An jedem Ende trägt die Querstange ein Kügelchen aus Platin. Offensichtlich ist das Trägheitsmoment dieses Pendels in Bezug auf die Schwingungsachse am kleinsten, wenn die Querstange parallel zu dieser Achse steht, am größten hingegen, wenn beide einen rechten Winkel miteinander bilden. Im ersten Fall schwingt das Pendel daher am schnellsten, im anderen am langsamsten. Bei geeignet gewählten Abmessungen des Apparats ist dadurch eine sehr feine Regulierung möglich, mit dem beachtlichen Vorteil, dass man den Gang korrigieren kann, ohne die Uhr anzuhalten. --- Über mehrere unter den Sternen \textsc{Flamstead}s gefundene Beobachtungen des Planeten Uranus, von Herrn \textsc{Burckhardt}. Ganz zufällig fand Herr \textsc{Burckhardt} bei anderen Nachforschungen in \textsc{Flamstead}s Historia Coelestis noch \emph{fünf} Uranusbeobachtungen (1712 März 22; 1715 Feb. 21. 22. 25 und April 18, alles nach altem Stil). Dieser Fund ist umso unerwarteter, als alle Astronomen den Gegenstand durch Herrn \textsc{Bode}s Nachforschung für ausgeschöpft hielten, und umso wertvoller, als die Beobachtungen von 1715 eine sehr gute Bestimmung der Opposition liefern, während die einzelne Beobachtung von 1690 noch mit einer Unsicherheit behaftet sein kann. \textsc{Flamstead} hat also den Planeten unwissentlich nicht weniger als sechsmal beobachtet, und die Entdeckung hätte ihm selbst gar nicht entgehen können, wenn er die Gewohnheit gehabt hätte, seine Beobachtungen immer gleich an Ort und Stelle zu überprüfen, wenn auch nur oberflächlich, und sie miteinander zu vergleichen. --- Über die kleinen Gleichungen in der Jupiterbewegung, von \emph{ebendemselben}. Herr \textsc{Burckhardt} gibt die numerischen Werte mehrerer von \textsc{Laplace} übergangener Gleichungen der \SourceLeaf{616}zweiten bis sechsten Ordnung an, von denen einige immerhin eine gewisse Bedeutung haben. Da die entsprechenden in der Theorie der Saturnbewegung übergangenen Gleichungen noch beträchtlich größer sein müssen, wäre es wohl möglich, dass dieser Umstand zu dem auffallenden Unterschied beigetragen hat, der sich bekanntlich bei der Bestimmung der Jupitermasse aus der Saturnbewegung und aus der Pallasbewegung ergeben hat. --- Von \emph{demselben}, Bemerkungen über verschiedene Fixsterne. --- Kometenbeobachtungen an der Königlichen Sternwarte in Paris von 1811---1815 (größtenteils schon anderswo gedruckt), nebst Elementen, die Herr \textsc{Nicollet} daraus abgeleitet hat. --- Anwendung der Wahrscheinlichkeitsrechnung auf die geodätischen Arbeiten, von \textsc{Laplace}. Die Wahrscheinlichkeitsrechnung findet dort Anwendung, wo mehr als das unbedingt Nötige beobachtet worden ist, sofern die Beobachtungen keine absolute Genauigkeit haben. Sie lehrt, wie man die Beobachtungen am vorteilhaftesten kombiniert und die relative Zuverlässigkeit der Ergebnisse bestimmt. Herr \textsc{Laplace} betrachtet hier den Fall, dass auf der Oberfläche einer Kugel zwei Punkte durch ein Dreiecksnetz verbunden sind, in den Dreiecken alle Winkel beobachtet und \emph{zwei} Grundlinien gemessen worden sind, eine am Anfang, die andere am Ende des Netzes. Die Winkelmessungen werden als fehlerbehaftet, die Grundlinie als fehlerfrei betrachtet. Gesucht werden die Länge des Bogens, der die beiden Punkte verbindet, und die Winkel, die er mit der ersten und der letzten Dreiecksseite bildet. Man sieht, dass dies nur eine sehr spezielle Aufgabe ist und die Anwendung der Wahrscheinlichkeitstheorie auf die höhere Geodäsie dadurch keineswegs ausgeschöpft ist.
\par\clearpage\endgroup
